Face recognition compares images by turning each into a vector and measuring distances. Two recipes produce the vectors: a classical linear subspace fitted to the pixels (Eigenfaces), and a deep network trained on millions of faces. This paper reports a controlled replication and measurement study of how much the second recipe buys, and where its claimed robustness to illumination comes from. On the 400-image Olivetti Faces dataset we compare a raw-pixel Eigenfaces map ($4096$-d) with a frozen VGGFace2-pretrained Inception-ResNet-v1 embedding ($512$-d, unit norm), using identical PCA settings for visualisation and computing every metric in full dimensionality. We add an open-set verification analysis under one protocol for both representations, and a webcam field study of one person under six lighting conditions.

The embedding wins on every identity metric: silhouette 0.106 to 0.497, 1-NN accuracy 94.5\% to 100\%, pair ROC area 0.881 to 0.9995, and open-set equal error rate 8.7\% to 0.22\% (at the protocol's 0.5\% false-reject resolution). At a matched 1\% false-reject rate, Map A accepts 48.6\% of impostors against 0.044\% for Map B. Two methodological findings matter as much. The raw-pixel 2-D plot preserves \emph{more} of its space's pairwise geometry (rank correlation 0.807) than the embedding's (0.517), so the prettier picture belongs to the worse space. And ``robust to lighting'' conflates two effects: brightness changes cost nothing beyond repeat-scan noise, while light \emph{direction} costs 3.83 times the frame-to-frame noise floor, with under-lighting worst. The work is a controlled replication, not a new method: Map A is fitted transductively on all 400 images and the deep model's training corpus is far larger than the test set, so the numbers measure a gap under favourable conditions.
